\chapter{Model Explainability, Visual Analytics \& Banking AI Governance}

\begin{tcolorbox}[colback=primarylight, colframe=primarydark, title={\textbf{What You Will Learn in this Chapter}}]
\begin{itemize}
    \item Regulatory imperatives for Explainable AI (XAI) in banking (FCRA Adverse Action, GDPR Art.~22, RBI guidelines)
    \item The Black-Box dilemma: Balancing complex non-linear models against legal interpretability
    \item Cooperative Game Theory foundations: Shapley values and axiomatic fairness
    \item SHAP (SHapley Additive exPlanations) and TreeSHAP algorithm mechanics
    \item Local Interpretable Model-agnostic Explanations (LIME)
    \item Visual forensic analytics: Sankey money flow diagrams, network graphs, and interactive dashboards
    \item Complete Python implementation of SHAP feature attribution on a loan denial decision
\end{itemize}
\end{tcolorbox}

\section{The Regulatory Imperative for Explainable AI (XAI)}

In non-regulated consumer domains (e.g., video recommendations), predictive accuracy is paramount and algorithmic obscurity is tolerated. In financial banking, however, algorithmic opacity is legally impermissible.

\begin{definitionbox}{Regulatory Mandates for Model Explainability}
\index{Adverse Action Notice}\index{Model Risk Management}
\begin{itemize}
    \item \textbf{Adverse Action Notices (FCRA / ECOA):} When a lender denies credit or raises interest rates based on an algorithmic score, the law requires delivery of specific, actionable reasons (e.g., ``debt-to-income too high'' or ``delinquency history'') to the consumer.
    \item \textbf{GDPR Article 22:} Prohibits solely automated individual decision-making having legal effects without meaningful human intervention and guarantees the consumer's ``Right to an Explanation.''
    \item \textbf{Model Risk Management (SR 11-7 / RBI Guidelines):} Central banks require institutions to rigorously validate that algorithmic decision engines do not rely on spurious correlations, discriminatory proxy features, or non-linear blind spots.
\end{itemize}
\end{definitionbox}

\begin{videocard}{IITM BS Lecture: Explainable AI}{https://www.youtube.com/watch?v=CQs-u6iu9Zo}
Conceptual foundations of algorithmic transparency, interpretability vs.~accuracy frontiers, and regulatory guardrails.
\end{videocard}

\section{Cooperative Game Theory: Shapley Values}

Developed by Nobel laureate Lloyd Shapley in 1953, \index{Shapley values}\textbf{Shapley values} calculate the fair attribution of total payout among a coalition of cooperative players based on their marginal contributions. In machine learning:
\begin{itemize}
    \item The ``game'' is the model's prediction task for a single transaction.
    \item The ``payout'' is the difference between the model's actual prediction and the base expected value.
    \item The ``players'' are the individual feature values ($x_1, x_2, \dots, x_M$) of the transaction.
\end{itemize}

\begin{formulabox}{The Shapley Value Formula}
\index{Shapley Value Formula}
The unique attribution $\phi_i$ assigned to feature $i$ is:
\begin{equation}
\phi_i(v) = \sum_{S \subseteq N \setminus \{i\}} \frac{|S|!(|N| - |S| - 1)!}{|N|!} \left[ v(S \cup \{i\}) - v(S) \right]
\end{equation}
where $N$ is the set of all features, $S$ is a subset of features excluding $i$, and $v(S)$ is the model outcome conditioned on feature subset $S$.
\end{formulabox}

\begin{theorembox}{The Four Axiomatic Properties of Shapley Values}
Shapley values represent the \textbf{only} attribution method satisfying all four fairness axioms:
\begin{enumerate}
    \item \textbf{Efficiency:} The sum of all feature attributions equals the total deviation from base prediction:
    \begin{equation}
    \sum_{i=1}^M \phi_i = f(\mathbf{x}) - E[f(X)]
    \end{equation}
    \item \textbf{Symmetry:} If features $i$ and $j$ contribute identically across all possible subsets ($v(S \cup \{i\}) = v(S \cup \{j\})$), then $\phi_i = \phi_j$.
    \item \textbf{Dummy Player:} If feature $i$ contributes nothing to any coalition ($v(S \cup \{i\}) = v(S)$), then $\phi_i = 0$.
    \item \textbf{Additivity:} For independent ensemble components ($f = f_1 + f_2$), the combined attribution is $\phi_i(f_1 + f_2) = \phi_i(f_1) + \phi_i(f_2)$.
\end{enumerate}
\end{theorembox}

\begin{videocard}{Masterclass Video: SHAP (SHapley Additive exPlanations) -- StatQuest}{https://www.youtube.com/watch?v=VB9Q6l0_2Ak}
Visual walk-through of coalitional games, marginal feature additions, and interpreting SHAP waterfall plots.
\end{videocard}

\begin{videocard}{IITM BS Lecture: SHAP Hands-On Tutorial}{https://www.youtube.com/watch?v=oK_Dex7jVEg}
Practical Python implementation using the `shap` library to generate force plots, summary bar charts, and dependence plots.
\end{videocard}

\begin{videocard}{IITM BS Lecture: Leveraging AI/ML for Financial Forensics}{https://www.youtube.com/watch?v=OpPJcD6qazA}
Synthesizing end-to-end artificial intelligence workflows into bank compliance architectures, fraud desks, and regulatory audits.
\end{videocard}

\begin{intuitionbox}{The Basketball Team Attribution}
Imagine a basketball team wins a game by 20 points over their average scoring margin. Did player A score those points alone, or did player B's defensive rebounds allow player A to take open shots? Shapley values simulate every possible lineup permutation (with and without Player A, with and without Player B) to calculate the exact, unarguable contribution of every teammate.
\end{intuitionbox}

\section{Visual Analytics in Financial Forensics}

Forensic reports presented to boards of directors, prosecutors, or judges must convey complex multi-hop evidence clearly through visual analytics:
\begin{itemize}
    \item \textbf{Sankey Diagrams:} Visualize volumetric conservation of money as illicit funds flow from a victim enterprise, pass through layering accounts, and aggregate in offshore tax shelters.
    \item \textbf{Network Graphs:} Expose hidden syndicates by plotting accounts, IP addresses, physical addresses, and directors as nodes, with transactions or corporate registry filings as connecting edges.
    \item \textbf{SHAP Waterfall Plots:} Provide auditable single-decision justifications for why a specific credit card transaction was blocked or a loan application was rejected.
\end{itemize}

\section{Data Science Relevance: SHAP Model Attribution Pipeline}

\begin{lstlisting}[caption={Automated Model Explainability with SHAP in Python}]
import numpy as np
import pandas as pd
from sklearn.ensemble import GradientBoostingClassifier

# 1. Simulate credit risk training dataset
np.random.seed(42)
n = 1500
income = np.random.uniform(20, 120, n)
debt_to_income = np.random.uniform(0.1, 0.8, n)
missed_payments = np.random.randint(0, 5, n)
credit_history_yrs = np.random.uniform(1, 20, n)

# Default risk latent score
risk = 2.5 * debt_to_income + 0.8 * missed_payments - 0.02 * income - 0.05 * credit_history_yrs
p_default = 1.0 / (1.0 + np.exp(-risk))
y_default = (np.random.rand(n) < p_default).astype(int)

df = pd.DataFrame({
    'Income': income,
    'DTI': debt_to_income,
    'Missed_Payments': missed_payments,
    'History_Years': credit_history_yrs
})

# 2. Train Gradient Boosted Decision Tree
model = GradientBoostingClassifier(n_estimators=60, max_depth=3, random_state=42)
model.fit(df, y_default)

# 3. Simulate marginal feature attribution (simplified Shapley calculation)
# For a specific rejected applicant:
applicant = pd.Series({'Income': 25.0, 'DTI': 0.72, 'Missed_Payments': 3, 'History_Years': 2.0})

base_prediction = model.predict_proba(df)[:, 1].mean()
applicant_pred = model.predict_proba(applicant.values.reshape(1, -1))[0, 1]

# Compute feature contribution differences
attributions = {}
for col in df.columns:
    # Marginal ablation: replace applicant value with baseline median
    ablated = applicant.copy()
    ablated[col] = df[col].median()
    p_ablated = model.predict_proba(ablated.values.reshape(1, -1))[0, 1]
    attributions[col] = round(applicant_pred - p_ablated, 4)

print(f"Base Portfolio Default Risk:  {base_prediction:.3f}")
print(f"Applicant Default Risk:       {applicant_pred:.3f}")
print("\nAdverse Action Explanations (Feature Attributions):")
for feat, val in sorted(attributions.items(), key=lambda x: x[1], reverse=True):
    print(f" - {feat:16s}: contribution = +{val:.4f} to risk")
\end{lstlisting}

\section{Forensic Visual Analytics: Building Dashboards in Tableau}

Interactive visual analytics bridge the communication divide between data science teams, forensic auditors, and courtroom judges:

\begin{tcolorbox}[colback=tealbg, colframe=tealborder, title={\textbf{Tableau Forensic Dashboard Blueprint}}]
\index{Tableau in Forensics}
\begin{enumerate}
    \item \textbf{Benford Conformity Chart with Parameter Thresholds:}
    \begin{itemize}
        \item Create a calculated field \texttt{Leading\_Digit} = \texttt{INT(LEFT(STR(ABS([Amount])), 1))}.
        \item Create calculated field \texttt{Theoretical\_Benford} = \texttt{LOG10(1 + 1/[Leading\_Digit]) * TOTAL(COUNT([Amount]))}.
        \item Plot Dual-Axis Bar (Observed count) vs.~Line (Expected Benford count), with color alert calculated field \texttt{ABS([Observed] - [Expected]) > [Tolerance\_Param]}.
    \end{itemize}
    \item \textbf{Transaction Velocity Scatter Matrix:}
    \begin{itemize}
        \item Set Columns = \texttt{DATEPART('hour', [Timestamp])} (0 to 23 hours), Rows = \texttt{[Amount]}.
        \item Color mark by \texttt{[Anomaly\_Score]} (from ML model export) and size mark by \texttt{[Daily\_Velocity\_Count]}.
        \item Enables instant visual isolation of high-value, high-velocity transactions executed during midnight hours.
    \end{itemize}
    \item \textbf{High-Risk Merchant Geographic Heatmap:}
    \begin{itemize}
        \item Drag \texttt{Latitude} / \texttt{Longitude} to map view, density colored by \texttt{SUM([Chargeback\_Count])}.
        \item Add interactive dashboard filter actions for Merchant Category Codes (MCC) to isolate merchant collusion clusters.
    \end{itemize}
\end{enumerate}
\end{tcolorbox}

\section*{Supplementary Lecture Videos}
\begin{itemize}
    \item \href{https://www.youtube.com/watch?v=CQs-u6iu9Zo}{W12\_L1: Explainable AI in Regulated Financial Domains}
    \item \href{https://www.youtube.com/watch?v=oK_Dex7jVEg}{W12\_L2: SHAP Hands-On Python Tutorial (TreeSHAP \& Waterfalls)}
    \item \href{https://www.youtube.com/watch?v=OpPJcD6qazA}{W12\_L3: Leveraging AI/ML for Enterprise Financial Forensics}
    \item \href{https://www.youtube.com/watch?v=VB9Q6l0_2Ak}{External Masterclass: SHAP: Model Interpretability (StatQuest)}
\end{itemize}

\newpage
\section{Practice Exercises}
\textit{Detailed step-by-step solutions are available in the Appendix.}

\begin{enumerate}
    \item Explain why the Fair Credit Reporting Act (FCRA) legally mandates adverse action disclosures. What legal vulnerability does a financial institution incur if it deploys an uninterpretable deep neural network to reject mortgage applicants?
    
    \item State and mathematically define the four foundational axioms of cooperative game theory that uniquely define Shapley values: Efficiency, Symmetry, Dummy Player, and Additivity.
    
    \item Contrast the algorithmic approaches of SHAP (SHapley Additive exPlanations) and LIME (Local Interpretable Model-agnostic Explanations). What theoretical advantages does TreeSHAP offer over LIME when explaining Gradient Boosted Decision Trees?
    
    \item Consider a simple cooperative game with 3 features ($x_1, x_2, x_3$). The model produces predictions for feature subsets as follows:
    $v(\emptyset) = 0.20$, $v(\{1\}) = 0.35$, $v(\{2\}) = 0.20$, $v(\{3\}) = 0.30$,
    $v(\{1, 2\}) = 0.50$, $v(\{1, 3\}) = 0.60$, $v(\{2, 3\}) = 0.30$, and $v(\{1, 2, 3\}) = 0.85$.
    Using the Shapley formula, compute the exact attribution $\phi_1$ for feature $1$.
    
    \item Describe how a forensic visual analytics dashboard combines Sankey money-flow diagrams with network community detection algorithms to trace funds siphoned across multiple nested corporate shell entities.
\end{enumerate}
